\clearpage
\SourceLeaf{264}
\begin{center}
{\large SECTION IV.}\par\medskip
\pdfbookmark[1]{On determining orbits while taking perturbations into account.}{liber-ii-sectio-iv}
\textit{On determining orbits while taking perturbations into account.}
\end{center}
\Needspace{6\baselineskip}
\begin{center}\hypertarget{native-264-article190}{}\pdfbookmark[2]{190.}{article-190}190.\end{center}
The perturbations that the motions of planets undergo through the action of
the other planets are so small and slow that they become perceptible only
after a longer interval. Over a shorter time --- or even over one or more
complete revolutions, depending on the circumstances --- the motion will
differ so little from motion along a perfect ellipse traced exactly according
to \textsc{Kepler}'s laws that the observations cannot reveal the deviation.
As long as this is the case, it would hardly be worth calculating the
perturbations prematurely. It suffices to fit a conic section, an osculating
one as it were, to the observations. Once the planet has been accurately
observed for a longer time, however, the effect of the perturbations will
eventually become apparent to such an extent that purely elliptic motion
can no longer reconcile all the observations exactly. Complete and lasting
agreement can then be obtained only by duly combining the perturbations
with the elliptic motion.

Determining the elliptic elements that must be combined with the perturbations
to represent the observations accurately presupposes knowledge of those
perturbations. Conversely, an accurate theory of perturbations requires
that the elements already be known approximately. The nature of the problem
precludes a perfectly accurate solution of this difficult task at the first
attempt. The highest degree of precision in the perturbations and the elements
can be reached only by repeatedly correcting them in turn. The first theory
of perturbations will therefore be based on the purely elliptic
\SourceLeaf{265}elements fitted approximately to the observations. We then
seek a new orbit that, when combined with these perturbations, fits the
observations as closely as possible. If it differs considerably from the
previous orbit, the perturbations must be worked out again on this basis.
These alternating corrections are repeated until the observations, the
elements, and the perturbations agree as closely as possible.
